\section[Proof of Theorem 6.2]{Proof of Theorem~\ref{thm-diophantine}}
\label{app-rings}

Consider the rings $\Z$ and $\D$, together with their respective
extensions $\Z[\sqrt2]$, $\Z[\omega]$ and $\D[\sqrt2]$, $\D[\omega]$,
as introduced in Section~\ref{sec-algebra}. We wish to give an
efficient method for solving equations of the form $t\da t=\xi$, for
given $\xi\in\D[\sqrt2]$ and unknown $t\in\D[\omega]$. To do this, we
first classify the primes in $\Z[\sqrt2]$ and $\Z[\omega]$, then show
how to solve the equation $t\da t=\xi$ in the case where
$\xi\in\Z[\sqrt2]$ and $t\in\Z[\omega]$, which we finally extend to
the general case. None of the results in this section are original;
they are well-known from the theory of cyclotomic fields, which goes
back to the 19th century work of Gauss and Kummer.  Nevertheless, we
hope that the following detailed treatment may be useful to readers
who are not experts in algebraic number theory. Moreover, we hope to
give sufficient details to enable an interested reader, in principle,
to implement the algorithm. This will also aid in our complexity
analysis.

A fundamental property of the rings $\Z$, $\Z[\sqrt2]$, and
$\Z[\omega]$ is that they are {\em Euclidean domains}; this implies
that the notions of divisibility, greatest common divisor, and unique
prime factorization all make sense in these rings.  Recall that in a
ring, a {\em unit} is an invertible element.  In a Euclidean domain,
we write $x\divides y$ if $x$ is a divisor of $y$, and $x\sim y$ if
$x\divides y$ and $y\divides x$; equivalently, $x\sim y$ iff there
exists a unit $u$ that $x u=y$. An element $x$ is {\em prime} if $x$
is not a unit, and $x=ab$ implies that either $a$ or $b$ is a
unit. Note that if $x$ is prime and $x\divides ab$, then $x\divides a$
or $x\divides b$; this follows from Euclid's algorithm.

\subsection[{Units in Z[sqrt 2]}]{Units in $\Z[\sqrt2]$}

\begin{definition}
  We say that $\xi\in\Z[\sqrt 2]$ is {\em doubly positive} if $\xi\geq
  0$ and $\xi\bul \geq 0$.
\end{definition}

\begin{lemma}\label{lem-units}
  The units of $\Z[\sqrt2]$ are of the form $u=(-1)^n\lambda^m$, where
  $\lambda=1+\sqrt 2$. Moreover, a unit $u$ is doubly positive if and
  only if $u$ is a square in $\Z[\sqrt2]$.
\end{lemma}

\begin{proof}
  Lemma~10 of {\cite{Selinger-newsynth}}.
\end{proof}

\begin{lemma}
  Let $\xi\in\Z[\sqrt2]$, and consider $n=\xi\bul\xi$. If $n$ is a unit
  in $\Z$, then $\xi$ is a unit in $\Z[\sqrt2]$.
\end{lemma}

\begin{proof}
  If $n$ is a unit, then there exists $m\in\Z$ such that $nm=1$, hence
  $\xi\xi\bul m=1$, hence $\xi$ is invertible in $\Z[\sqrt2]$.
\end{proof}

%----------------------------------------------------------------------
\subsection[{Primes in Z[sqrt 2]}]{Primes in $\Z[\sqrt2]$}

\begin{lemma}\label{lem-prime-criterion}
  Let $\xi\in\Z[\sqrt2]$, and consider $n=\xi\bul\xi$. If $n$ is prime
  in $\Z$, then $\xi$ is prime in $\Z[\sqrt2]$.
\end{lemma}

\begin{proof}
  Suppose $\xi=\alpha\beta$ in $\Z[\sqrt 2]$, and consider $n =
  \xi\bul\xi = \alpha\bul\alpha \beta\bul\beta$. Since
  $\alpha\bul\alpha$ and $\beta\bul\beta$ are integers and $n$ is
  prime, we must have that either $\alpha\bul\alpha$ or
  $\beta\bul\beta$ is a unit in $\Z$; hence $\alpha$ or $\beta$ is a
  unit in $\Z[\sqrt 2]$. So $\xi$ is prime.
\end{proof}

\begin{lemma}
  For every prime $\xi$ of $\Z[\sqrt2]$, there exists a unique (up to
  a unit) prime $p$ of $\Z$ such that $\xi \divides p$.
\end{lemma}

\begin{proof}
  To show existence, consider $n=\xi\bul\xi$. Note that $\xi\neq 0$,
  hence $n\neq 0$. Let $n=p_1p_2\cdots p_k$ be a prime factorization
  of $n$. Since $\xi \divides p_1p_2\cdots p_k$ and $\xi$ is prime in
  $\Z[\sqrt2]$, there exists some $i$ such that $\xi\divides p_i$.

  To show uniqueness, assume $\xi\divides p$ and $\xi\divides q$, where $p\not\sim q$.
  Then $\gcd(p,q)=1$, hence by Euclid's algorithm, we can write
  $1=np+mq$, for integers $n$ and $m$. Then $\xi\divides np+mq=1$, which
  is absurd since $\xi$ is prime. 
\end{proof}

\begin{lemma}\label{lem-z2-cases}
  Let $\xi$ be a prime of $\Z[\sqrt2]$, and let $p$ be the unique (up
  to a unit) prime of $\Z$ such that $\xi\divides p$. Then exactly one of
  the following holds: $\xi\sim p$ or $\xi\bul\xi\sim p$.
\end{lemma}

\begin{proof}
  First note that at most one of these properties can hold, because
  otherwise $\xi\sim\xi\bul\xi$, which implies that $\xi\bul$ is a
  unit, which is absurd since it is prime.

  We now show that at least one of the properties holds. Since
  $\xi\divides p$ and $p$ is an integer, we have $\xi\bul\divides p$, hence
  $\xi\bul\xi\divides p^2$. Also, $\xi\bul\xi$ is an integer, so either
  $\xi\bul\xi\sim 1$, $\xi\bul\xi\sim p$, or $\xi\bul\xi\sim p^2$. The
  first of these cases is absurd, since $\xi$ would then be a unit. In
  the second case, we have $\xi\bul\xi\sim p$, which is to be
  shown. In the third case, since $\xi\divides p$, there is
  $\alpha\in\Z[\sqrt2]$ such that $\xi\alpha=p$. Then we have
  $\xi\bul\xi\alpha\bul\alpha=p^2\sim\xi\bul\xi$, which implies that
  $\alpha$ is a unit, hence $\xi\sim p$.
\end{proof}

\begin{lemma}\label{lem-prime-factor-z2}
  Let $p$ be a prime of $\Z$. Then the prime factorization of $p$ in
  $\Z[\sqrt2]$ consists of one or two factors. 
\end{lemma}

\begin{proof}
  Let $\xi$ be some prime factor of $p$ in $\Z[\sqrt2]$. By
  Lemma~\ref{lem-z2-cases}, either $\xi\sim p$ or $\xi\bul\xi\sim
  p$. Either way, this gives a prime factorization of $p$ in $\Z[\sqrt2]$.
\end{proof}

We now determine the prime factorization in $\Z[\sqrt2]$ of every
prime $p$ of $\Z$. It turns out that there are 5 cases, depending
whether $p$ is even, or $p\equiv 1, 3, 5, 7\mmod{8}$.

\begin{lemma}\label{lem-2}
  The prime factorization of $p=2$ in $\Z[\sqrt2]$ is $p =
  \sqrt{2}\cdot \sqrt{2}$. The prime factorization of $p=-2$ in
  $\Z[\sqrt2]$ is $p=-\sqrt{2}\cdot\sqrt{2}$.
\end{lemma}

\begin{proof}
  We only need to show that $\sqrt2$ is prime in $\Z[\sqrt2]$, but
  this follows from Lemma~\ref{lem-prime-criterion}.
\end{proof}

\begin{lemma}\label{lem-3-5}
  Let $p$ be a prime of $\Z$ such that $p\equiv 3\mmod{8}$ or $p\equiv
  5\mmod{8}$. Then $p$ is prime in $\Z[\sqrt2]$.
\end{lemma}

\begin{proof}
  Let $\xi$ be a prime factor of $p$ in $\Z[\sqrt2]$. By
  Lemma~\ref{lem-z2-cases}, we know that $\xi\sim p$ or
  $\xi\bul\xi\sim p$. In the former case, $p$ is prime and we are
  done. In the latter case, writing $\xi=a+b\sqrt2$, we have
  $\xi\bul\xi = a^2-2b^2 =\pm p$, hence $a^2-2b^2\equiv\pm3\mmod{8}$.
  By easy case distinction, we see that $a^2$ can only be $0$, $1$, or
  $4$ modulo $8$, and $2b^2$ can only be $0$ or $2$ modulo $8$, so
  $a^2-2b^2\equiv\pm3\mmod{8}$ is plainly impossible.
\end{proof}

\begin{lemma}\label{lem-associate-primes}
  Let $\xi\in\Z[\sqrt2]$ such that $\xi\sim\xi\bul$. Then either
  $\xi\sim n$, or $\xi\sim n\sqrt{2}$, for some $n\in\Z$.
\end{lemma}

\begin{proof}
  By assumption, $\xi\sim\xi\bul$, so there exists a unit
  $u\in\Z[\sqrt2]$ such that $\xi\bul=u\xi$. By Lemma~\ref{lem-units}, the 
  units of $\Z[\sqrt2]$ are exactly of the form $u=(-1)^n\lambda^m$, where
  $\lambda=1+\sqrt2$. Applying the automorphism to $\xi\bul=u\xi$, we get 
  $\xi=u\bul\xi\bul$, hence $\xi\bul\xi = u\bul u\xi\bul\xi$, hence 
  $u\bul u=1$, hence $(\lambda\bul\lambda)^m = 1$. But $\lambda\bul\lambda=-1$, 
  so $m$ is even; say $m=2k$. Let $\xi' = \lambda^k\xi$. Note that
  $\xi'\sim\xi$. Then we have
  \begin{equation}\label{eqn-xi-prime}
    \xi'^\bullet = (\lambda\bul)^k\xi\bul = (-\lambda)^{-k} u\xi
    = (-1)^{-k}\lambda^{-k}(-1)^n\lambda^{m}\xi = \pm\lambda^k\xi = \pm \xi'.
  \end{equation}
  We can write $\xi' = a+b\sqrt2$ for some $a,b\in\Z$. From
  {\eqref{eqn-xi-prime}}, we have either $a-b\sqrt 2 = a+b\sqrt 2$, or
  $a-b\sqrt 2 = -(a+b\sqrt 2)$. In the first case, $b=0$, and therefore
  $\xi\sim\xi'= a$. In the second case, $a=0$, and therefore
  $\xi\sim\xi'= b\sqrt 2$.
\end{proof}

\begin{lemma}\label{lem-1-7}
  Let $p$ be a prime of $\Z$ such that $p\equiv 1\mmod{8}$ or $p\equiv
  7\mmod{8}$. Then $p$ has a prime factorization of the form $p\sim
  \xi\bul\xi$ in $\Z[\sqrt2]$; moreover, $\xi\not\sim\xi\bul$, so that
  the two prime factors are distinct.
\end{lemma}

\begin{proof}
  It is well-known (by quadratic reciprocity) that if $p$ is a prime
  with $p\equiv\pm1\mmod{8}$, then $2$ is a quadratic residue modulo
  $p$. Therefore, the equation $x^2\equiv 2\mmod{p}$ has an integer
  solution; let $x$ be such a solution. Let $\alpha=x+\sqrt2$. Then
  $\alpha\bul\alpha=x^2-2\equiv0\mmod{p}$, and hence
  $p\divides\alpha\bul\alpha$. On the other hand, clearly $p\nmid\alpha$
  (since $\alpha/p\not\in\Z[\sqrt2]$), and similarly
  $p\nmid\alpha\bul$, which shows that $p$ is not a prime in
  $\Z[\sqrt2]$. By Lemma~\ref{lem-z2-cases}, $p\sim\xi\bul\xi$ for
  some prime $\xi$ of $\Z[\sqrt 2]$. 

  The final claim follows from Lemma~\ref{lem-associate-primes},
  because if $\xi\sim\xi\bul$, then either $\xi\sim n$ or $\xi\sim
  n\sqrt2$. In the first case, $n^2\divides \xi\bul\xi \sim p$, but $p$ is
  prime, so that $n=\pm 1$, contradicting that $\xi$ is prime. In the
  second case, $2n^2\divides\xi\bul\xi\sim p$, contradicting the
  assumption that $p$ is odd.
\end{proof}

\begin{lemma}\label{lem-prime-factorization-z2}
  Let $p$ be a prime of $\Z$. A prime factorization of $p$ in
  $\Z[\sqrt2]$ can be computed in probabilistic polynomial time.
\end{lemma}

\begin{proof}
  Assume without loss of generality that $p>0$. If $p=2$, then
  $p=\rt{2}$ is the desired prime factorization by
  Lemma~\ref{lem-2}. If $p\equiv 3,5\mmod{8}$, then $p$ is already
  prime in $\Z[\sqrt2]$ by Lemma~\ref{lem-3-5}. The remaining case is
  when $p\equiv 1,7\mmod{8}$. In this case the prime factorization is
  of the form $p\sim\xi\bul\xi$ by Lemma~\ref{lem-1-7}. Moreover, the
  proof of Lemma~\ref{lem-1-7} indicates how such $\xi$ can be
  computed, namely as $\xi=\gcd(p,x+\sqrt2)$, where $x$ is a solution
  of $x^2\equiv 2\mmod{p}$. The equation $x^2\equiv 2\mmod{p}$ can be
  solved in probabilistic polynomial time by a well-known algorithm,
  see {\cite{Rabin1980}}.
\end{proof}

\subsection[{Primes in Z[omega]}]{Primes in $\Z[\omega]$}

\begin{lemma}
  Let $\xi$ be a prime in $\Z[\sqrt2]$. Then either $\xi$ is prime in
  $\Z[\omega]$, or else $\xi\sim t\da t$ where $t$ is some prime of
  $\Z[\omega]$. In particular, the prime factorization
  of $\xi$ in $\Z[\omega]$ consists of either one or two factors. 
\end{lemma}

\begin{proof}
  This is similar to the proof of Lemma~\ref{lem-prime-factor-z2}.
  Let $t$ be some prime factor of $\xi$ in $\Z[\omega]$. Then
  $t\divides \xi$, therefore $t\da\divides \xi\da=\xi$, therefore
  $t\da t\divides \xi^2$ in $\Z[\omega]$. But both $t\da t$ and $\xi$
  are elements of $\Z[\sqrt2]$, so $t\da t\divides \xi^2$ in
  $\Z[\sqrt2]$. Since $\xi$ is prime in $\Z[\sqrt2]$, there are only
  three possibilities: $t\da t\sim 1$, $t\da t\sim \xi$, and $t\da
  t\sim \xi^2$. In the first case, $t$ is a unit, contradicting the
  fact that it is prime in $\Z[\omega]$. In the second case, we are
  done. In the third case, since $t\divides \xi$, there is
  $a\in\Z[\omega]$ such that $at=\xi$. Then $a\da at\da
  t=\xi\da\xi=\xi^2 = t\da t$, which implies $a\da a=1$. Therefore $a$
  is a unit, and $t\sim \xi$. Therefore $\xi$ is prime in $\Z[\omega]$
  and we are done.
\end{proof}

\subsection[{The Diophantine equation t* t = xi}]
{The Diophantine equation $t\da t=\xi$}

We are interested in solving equations of the form 
\begin{equation}\label{eqn-tdat}
t\da t=\xi,
\end{equation}
where $\xi\in\D[\sqrt2]$ is given, and $t\in\D[\omega]$ is unknown.

\begin{definition}
  Recall that for elements $\xi,\xi'\in\Z[\sqrt2]$, the notation
  $\xi\sim\xi'$ means that $\xi,\xi'$ differ by a unit, i.e., there
  exists a unit $u\in\Z[\sqrt2]$ such that $\xi=u\xi'$. We extend this
  notation also to the ring $\D[\sqrt2]$: for $\xi,\xi'\in\D[\sqrt2]$,
  we write $\xi\sim\xi'$ iff there exists a unit $u\in\Z[\sqrt2]$ such
  that $\xi=u\xi'$. Note that we have taken $u$ to be a unit of the
  ring $\Z[\sqrt2]$, not of $\D[\sqrt2]$.
\end{definition}

It will often be convenient to replace {\eqref{eqn-tdat}} by the
following weaker condition. 

\begin{definition}
  We say that $\xi\in\D[\sqrt2]$ is \emph{$\dagger$-decomposable} if
  the equation
  \begin{equation}\label{eqn-tdat-associates}
    t\da t\sim\xi
  \end{equation} 
  has a solution $t\in\D[\omega]$.
\end{definition}

Solutions to {\eqref{eqn-tdat}} can be recovered from solutions to
{\eqref{eqn-tdat-associates}} by using the following observation:

\begin{lemma}\label{lem-sim}
  Let $\xi\in\D[\sqrt2]$. Then equation {\eqref{eqn-tdat}} has a
  solution if and only if $\xi$ is doubly positive and
  $\dagger$-decomposable.
\end{lemma}

\begin{proof}
  If $\xi$ is a solution to {\eqref{eqn-tdat}}, then it is obviously
  $\dagger$-decomposable. It is also doubly positive by
  Lemma~\ref{lem-nec-cond}. Conversely, assume $t\da t\sim
  \xi$. Then there exists a unit $u$ of $\Z[\sqrt2]$ such that
  $\xi=ut\da t$. Since both $\xi$ and $t\da t$ are doubly positive, it
  follows that $u$ is doubly positive, and hence $u$ is a square of
  the ring $\Z[\sqrt2]$ by Lemma~\ref{lem-units}; say $u=v^2$. Since
  $v\in\Z[\sqrt2]$, we have $v=v\da$. Setting $t'=vt$, we have
  $\xi=v\da vt\da t = t'^{\dagger}t'$, which finishes the proof.
\end{proof}

%----------------------------------------------------------------------
\subsection[{The case xi in Z[root 2]}]
{The case $\xi\in\Z[\sqrt 2]$}

\begin{lemma}\label{lem-zomega-domega}
  Suppose $t\da t\sim \xi$ for $t\in\D[\omega]$ and $\xi\in\Z[\sqrt
  2]$. Then $t\in\Z[\omega]$. 
\end{lemma}

\begin{proof}
  Note that, in $\Z[\omega]$, we have $\rt{k}\divides t$ if and only
  if $2^k\divides t\da t$. Choose $t'\in\Z[\omega]$ and $k\geq 0$ such
  that $t=t'/\rt{k}$. Then $t'^{\dagger}t' = 2^k\xi$, hence
  $2^k\divides t'^{\dagger}t'$, hence $\rt{k}\divides t'$, hence
  $t\in\Z[\omega]$. 
\end{proof}

\begin{lemma}\label{lem-gcd}
  If $x,y,z$ are three elements of a Euclidean domain, then
  \[\gcd(xy,z)\divides \gcd(x,z)\cdot\gcd(y,z).
  \]
\end{lemma}

\begin{proof}
  By considering the prime factorization of $z$. 
\end{proof}

\begin{lemma}\label{lem-alpha-beta}
  Suppose $\xi=\alpha\beta$, where $\alpha,\beta\in\Z[\sqrt2]$ and
  $\gcd(\alpha,\beta)=1$. Then $\xi$ is $\dagger$-decomposable iff
  $\alpha$ and $\beta$ are $\dagger$-decomposable.
\end{lemma}

\begin{proof}
  For the right-to-left implication, assume $\alpha\sim s\da s$ and
  $\beta\sim r\da r$. Clearly $\xi=\alpha\beta\sim (sr)\da sr$. For
  the left-to-right implication, assume $\xi\sim t\da t$. Note that
  $t\in\Z[\omega]$ by Lemma~\ref{lem-zomega-domega}. To show that
  $\alpha$ is $\dagger$-decomposable, let $s=\gcd(t, \alpha)$.  We
  claim that $s\da s\sim\alpha$. Clearly, since $s\divides\alpha$ and
  $s\da\divides\alpha$, we have $s\da s\divides \alpha^2$. One the
  other hand, we know that $s\da s\divides t\da t \sim
  \xi=\alpha\beta$. Since $\alpha$ and $\beta$ are relatively prime,
  it follows that $s\da s\divides \alpha$. Conversely, note that by
  Lemma~\ref{lem-gcd}, $\alpha = \gcd(t\da t,\alpha) \divides
  \gcd(t\da,\alpha)\cdot \gcd(t,\alpha) = s\da s$. Therefore $s\da
  s\sim\alpha$ and $\alpha$ is $\dagger$-decomposable. The argument
  for $\beta$ is similar.
\end{proof}

\begin{lemma}\label{lem-1235}
  Suppose that $\xi\in\Z[\sqrt2]$ is prime. Let $p>0$ be the unique
  positive prime in $\Z$ such that $\xi\divides p$. Then $\xi$ is
  $\dagger$-decomposable if and only if $p=2$ or $p\equiv
  1,3,5\mmod{8}$.
\end{lemma}

\begin{proof}
  We consider each case in turn. If $p=2$, then $\xi\sim\sqrt2$. Let
  $\delta=1+\omega$; a simple calculation shows that $\delta\da
  \delta=\lambda\sqrt2\sim\xi$. So $\xi$ is $\dagger$-decomposable.

  If $p\equiv 1\mmod{4}$, then by quadratic reciprocity, there exists
  some integer $u$ such that $u^2\equiv -1\mmod{p}$. Therefore
  $\xi\divides p\divides u^2+1 = (u+i)(u-i)$. Let $t=\gcd(\xi,u+i)$.
  We claim that $\xi\sim t\da t$. Note that $t\divides \xi$, hence
  $t\da\divides\xi\da=\xi$, hence $t\da t\divides \xi^2$. Since $\xi$
  is prime in $\Z[\sqrt2]$, there are 3 possibilities: $t\da t\sim 1$,
  $t\da t\sim \xi$, or $t\da t\sim \xi^2$. The first case is not
  possible, because in this case, $t$ would be a unit, so that $\xi$
  is relatively prime to $u+i$, hence to $u-i$, hence to $u^2+1$,
  contradicting $\xi\divides u^2+1$. In the second case, we have $t\da
  t\sim \xi$, which was to be shown. In the third case, we have $t\da
  t\sim\xi^2$. Since $t\divides \xi$, there exists some
  $s\in\Z[\omega]$ such that $ts=\xi$. But then $t\da ts\da s =
  \xi\da\xi=\xi^2\sim t\da t$, so that $s$ is a unit. In this case, we
  have $t\sim \xi$, therefore $\xi\divides u+i$, therefore also
  $\xi=\xi\da\divides u-i$, hence $\xi\divides (u+i)-(u-i)\sim 2$,
  contradicting the fact that $p$ is the only prime integer divisible
  by $\xi$.

  The case $p\equiv 3\mmod{8}$ is very similar. In this case, by
  quadratic reciprocity, $-2$ is a square modulo $p$, so that there
  exists some integer $u$ such that $u^2\equiv-2\mmod{p}$. Therefore
  $\xi\divides p\divides u^2+2 = (u+i\sqrt2)(u-i\sqrt2)$. Let
  $t=\gcd(\xi,u+i\sqrt2)$.  We claim that $\xi\sim t\da t$. Note that
  $t\divides\xi$, hence $t\da\divides\xi\da=\xi$, so $t\da t\divides
  \xi^2$. So again we have the three possibilities $t\da t\sim 1$,
  $t\da t\sim \xi$, or $t\da t\sim \xi^2$. The first case is not
  possible, because in this case, $t$ would be a unit, so that $\xi$
  is relatively prime to $u+i\sqrt2$, hence to $u-i\sqrt2$, hence to
  $u^2+2$, contradicting $\xi\divides u^2+2$. In the second case, we
  have $t\da t\sim \xi$, which was to be shown. In the third case, we
  have $t\da t\sim\xi^2$. Since $t\divides \xi$, there exists some
  $s\in\Z[\omega]$ such that $ts=\xi$. But then $t\da ts\da s =
  \xi\da\xi=\xi^2\sim t\da t$, so that $s$ is a unit. In this case, we
  have $t\sim \xi$, therefore $\xi\divides u+i\sqrt2$, therefore also
  $\xi=\xi\da\divides u-i\sqrt2$, hence $\xi\divides
  i(u-i\sqrt2)-i(u+i\sqrt2)=2\sqrt2$. On the other hand, $\xi\divides
  p$, therefore $\xi\divides\gcd(p,2\sqrt2)=1$, contradicting the fact
  that $\xi$ is not a unit.

  Finally, if $p\equiv 7\mmod{8}$, then $p\sim\xi\bul\xi$ by
  Lemma~\ref{lem-1-7}. Assume that $\xi$ is $\dagger$-decomposable as
  $\xi\sim t\da t$. Then
  \[ (t\da t)\bul(t\da t) \sim \xi\bul\xi\sim p.
  \]
  Note that $t\bul t\in\Z[i]$, so we can write $t\bul t=a+bi$ for some
  $a,b\in\Z$. But then we have
  \[ p \sim (t\bul t)(t\bul t)\da = (a+bi)(a-bi) = a^2+b^2.
  \]
  But $a^2$ and $b^2$ can only be congruent to $0$, $1$, or $4$ modulo
  8, contradicting $p\equiv 7\mmod{8}$. Therefore $\xi$ is not
  $\dagger$-decomposable, which is what was claimed.
\end{proof}

\begin{lemma}\label{lem-1235-powers}
  Suppose $\xi\in\Z[\sqrt2]$ is prime. Let $p>0$ be the unique
  positive prime in $\Z$ such that $\xi\divides p$. Let $m$ be a
  positive integer. Then $\xi^m$ is $\dagger$-decomposable if and only
  if $m$ is even or $p\equiv 1,2,3,5\mmod{8}$.
\end{lemma}

\begin{proof}
  If $m$ is even, then note that $\xi\in\Z[\sqrt2]$, so $\xi=\xi\da$;
  therefore, $t=\xi^{m/2}$ will be a solution. If $p\equiv
  1,2,3,5\mmod{8}$, then by Lemma~\ref{lem-1235}, there exists
  $s\in\Z[\omega]$ with $s\da s\sim\xi$; therefore $t=s^m$ is a
  solution of $t\da t\sim\xi^m$. The only remaining case is then $m$
  is odd and $p\equiv 7\mmod{8}$. In this case, there can be no
  solution. For assume on the contrary that $t\da t=\xi^m$. Then as in
  the proof of Lemma~\ref{lem-1235}, we can write $t\bul t=a+bi$ for
  some $a,b\in\Z$, and we get
  \[ p^m \sim (t\bul t)(t\bul t)\da = (a+bi)(a-bi) = a^2+b^2,
  \]
  contradicting $p^m\equiv 7\mmod{8}$.
\end{proof}

\begin{remark}\label{rem-1235-eff}
  The proofs of Lemmas~\ref{lem-1235} and {\ref{lem-1235-powers}} are
  constructive, and immediately yield efficient algorithms for
  determining $t$, provided that we have an efficient method of
  solving $u^2\equiv-1\mmod{p}$ when $p$ is a prime congruent to
  $1\mmod{4}$ and of solving $u^2\equiv-2\mmod{p}$ when $p$ is a prime
  congruent to $3\mmod{8}$. These last problems can be solved in
  probabilistic polynomial time by a well-known algorithm, see
  {\cite{Rabin1980}}.
\end{remark}

\begin{lemma}\label{lem-solve}
  Given $\xi\in\Z[\sqrt2]$, together with its prime factorization in
  $\Z[\sqrt2]$, there exists an algorithm that determines, in
  probabilistic polynomial time, whether the equation $t\da t\sim\xi$ has
  a solution or not, and finds a solution if there is one.
\end{lemma}

\begin{proof}
  Let $\xi\sim \xi_1^{m_1}\xi_2^{m_2}\cdots\xi_k^{m_k}$ be a prime
  factorization of $\xi$ in $\Z[\sqrt2]$, where $\xi_1,\ldots,\xi_k$
  are distinct (i.e., pairwise non-associate) primes. By
  Lemma~\ref{lem-alpha-beta}, $t\da t\sim\xi$ has a solution if and
  only if $t\da t\sim\xi_i^{m_i}$ has a solution for all $i$.  By
  Lemma~\ref{lem-1235-powers}, $t\da t\sim\xi_i^{m_i}$ has a solution
  if and only if $m_i$ is even or $p\equiv 1,2,3,5\mmod{8}$. Since
  these conditions are easy to check, we can therefore determine the
  existence of a solution efficiently, i.e., in polynomial time.

  Moreover, once a solution has been shown to exist, the actual
  solution can be efficiently computed. Namely, for each $i$, find
  $t_i$ such that $t_i\da t_i\sim\xi_i^{m_i}$ as in
  Lemma~\ref{lem-1235-powers}. Then $t\sim t_1t_2\cdots t_k$ satisfies
  $t\da t\sim \xi_1^{m_1}\xi_2^{m_2}\cdots\xi_k^{m_k}\sim\xi$. By
  Remark~\ref{rem-1235-eff}, all of this can be computed in
  probabilistic polynomial time.
\end{proof}

\begin{proposition}\label{prop-diophantine}
  Let $\xi\in\Z[\sqrt2]$, and let $n=\xi\bul\xi$. Note that $n$ is an
  integer. Given the prime factorization of $n$, there exists an
  algorithm that determines, in probabilistic polynomial time, whether
  the equation $t\da t\sim\xi$ has a solution or not, and finds a
  solution if there is one.
\end{proposition}

\begin{proof}
  Suppose that $n=\pm p_1^{m_1}\cdots p_k^{m_k}$ is the prime
  factorization of $n$, where $p_1,\ldots,p_k$ are distinct positive
  primes.  Each $p_i$ can be efficiently factored into primes of
  $\Z[\sqrt2]$ by Lemma~\ref{lem-prime-factorization-z2}; this yields
  the prime factorization of $n=\xi\bul\xi$ in $\Z[\sqrt2]$. From
  this, it is easy to obtain a prime factorization of $\xi$. The rest
  of the claim then follows from Lemma~\ref{lem-solve}.
\end{proof}

%----------------------------------------------------------------------
\subsection[{The case xi in D[root 2]}]
{The case $\xi\in\D[\sqrt 2]$}

The following lemma can be used to reduce the problem of
$\dagger$-decomposability in $\D[\sqrt2]$ to $\dagger$-decomposability
in $\Z[\sqrt2]$.

\begin{lemma}\label{lem-Z-to-D-new}
  Consider $\xi\in\D[\sqrt 2]$. Then $\xi$ is $\dagger$-decomposable
  if and only if $\sqrt2\,\xi$ is $\dagger$-decomposable.
\end{lemma}

\begin{proof}
  Recall that $\delta=1+\omega$ satisfies $\delta\da\delta=\lambda\sqrt2 \sim
  \sqrt2$. Also note that $\delta$ is invertible in $\D[\omega]$;
  specifically, $\delta\inv = \delta\lambda\inv\omega\inv/\sqrt2 \in\D[\omega]$.
  Assume that $\xi$ is $\dagger$-decomposable, so that $t\da
  t\sim\xi$. Letting $t'=\delta t$, we have
  $t'^{\dagger}t'=\delta\da\delta t\da t \sim \sqrt2 t\da t\sim
  \sqrt2\,\xi$, so that $\sqrt2\,\xi$ is $\dagger$-decomposable.
  The converse is proved similarly, using $\delta\inv$ instead of
  $\delta$. 
\end{proof}

We can now prove Theorem~\ref{thm-diophantine}, whose statement we 
reproduce here.

\begin{un-theorem}
  Let $\xi\in\D[\sqrt 2]$. Note that $\xi\bul\xi\in\D$, so we can
  write $\xi\bul\xi=\frac{n}{2^\ell}$ for some $n\in\Z$ and $\ell\in\N$. 
  There exists a probabilistic algorithm which, given $\xi$ and, in 
  case $n\neq 0$, a prime factorization of $n$, determines whether or 
  not the equation $t\da t =\xi$ has a solution, and finds a 
  solution if there is one. Moreover, the expected runtime of this 
  algorithm is polynomial in the size of $n$.
\end{un-theorem}

\begin{proof}
  Let $\xi$ be an element of $\D[\sqrt 2]$ with $\xi\bul\xi=
  \frac{n}{2^\ell}$ for some $n\in\Z$. First, the algorithm can easily
  check whether $\xi$ is doubly positive, and if it is not, there is
  no solution. Next, if $n=0$, then $\xi=0$, and $t=0$ is obviously a
  solution, so there is nothing to do. Otherwise, let $\xi'= {\rt
    \ell}\xi$. Since $\xi'\in \Z[\sqrt 2]$ and $\xi'^{\bullet} \xi'=n$, and
  a factorization of $n$ is given, we can use
  Proposition~\ref{prop-diophantine} to efficiently determine whether
  $s\da s\sim \xi'$ has a solution. By
  Lemma~\ref{lem-Z-to-D-new} this is the case if and only if the equation
  $t\da t\sim\xi$ also has a solution, in which case it is given by
  $t=\delta^{-\ell} s$. Finally, since $\xi$ is doubly positive, we can
  solve $t'^{\dagger}t'=\xi$ by Lemma~\ref{lem-sim}.
\end{proof}

We conclude this appendix with a useful fact: the algorithm of
Theorem~\ref{thm-diophantine} always succeeds in case $n$ is a prime
that is congruent to $1$ modulo $8$.

\begin{proposition}\label{prop-prime-1mod8}
  Let $\xi=\D[\sqrt 2]$ be doubly positive, with $\xi\bul\xi=
  \frac{n}{2^\ell}$ for some $n\in\Z$. If $n$ is prime and $n\equiv
  1\mmod{8}$, then the equation $t\da t=\xi$ has a solution.
\end{proposition}

\begin{proof}
  Let $\xi'={\rt \ell}\xi$. Then $\xi'\in\Z[\sqrt 2]$ and
  $\xi'^{\bullet}\xi'=n$. Then $\xi'$ is prime by
  Lemma~\ref{lem-prime-criterion} and $\dagger$-decomposable by
  Lemma~\ref{lem-1235}. By Lemma~\ref{lem-Z-to-D-new}, $\xi$ is
  $\dagger$-decomposable, so that the equation $t\da t=\xi$ can be
  solved by Lemma~\ref{lem-sim}.
\end{proof}
